\documentclass[main.tex]{subfiles}
\begin{document}

\chapter{Sufficiency and the Rao--Blackwell Theorem}\label{ch:yeterlilik}

A sample of $n$ observations contains many details that carry no information about $\theta$ (the order of the observations, a particular rearrangement, and so on). A sufficient statistic\index{sufficient statistic@sufficient statistic} is a summary that compresses the sample while retaining \emph{all} the information about $\theta$: once it is known, the rest of the raw data says nothing about $\theta$. The Rao--Blackwell\index{rao blackwell theorem@Rao--Blackwell theorem} theorem links this idea to estimation theory: any estimator can be improved, \emph{never made worse}, by conditioning on a sufficient statistic. This chapter is the first half of the UMVUE\index{umvue@UMVUE} theory of Chapter~\ref{ch:tamlik}.

\section{Definitions}

\begin{definition}[Sufficient Statistic]\label{def:yeterli}
Let $T=T(\X)$ be a statistic. If the conditional distribution of $\X$ given $T$ does not depend on $\theta$, i.e.\ if $\P_\theta(\X\in A\mid T=t)$ is free of $\theta$ for every $A$ and every $t$, then $T$ is called \emph{sufficient} for $\theta$.
\end{definition}

\begin{remark}[Interpretation]
If $T$ is sufficient, a statistician who knows $T=t$ can, without knowing $\theta$, generate an $\X^*$ from the known conditional distribution; $\X^*$ has the same distribution as $\X$. Hence every inference that can be made with $\X$ can be made with the same performance using $T$ alone. This is the rationale for the \emph{sufficiency principle}.
\end{remark}

\begin{definition}[Minimal Sufficient Statistic]
$T$ is called \emph{minimal sufficient\index{minimal sufficient statistic@minimal sufficient statistic}} if it is sufficient and, for every other sufficient statistic\index{sufficient statistic@sufficient statistic} $U$, there is a function $h$ with $T=h(U)$. A minimal sufficient statistic is the greatest data compression achievable without loss of information.
\end{definition}

\begin{definition}[Exponential Family]\label{def:ustel-aile}
A family of densities (or pmfs) is called a \emph{$k$-parameter exponential family\index{exponential family@exponential family}} if it can be written as
\[
f(x;\bm\theta)=h(x)\,c(\bm\theta)\exp\Big\{\sum_{j=1}^k w_j(\bm\theta)\,t_j(x)\Big\}.
\]
The form $f(x;\bm\eta)=h(x)\exp\{\sum\eta_jt_j(x)-A(\bm\eta)\}$ obtained through $\eta_j=w_j(\bm\theta)$ is called the \emph{natural (canonical) form}; $\bm\eta$ is the natural parameter\index{natural parameter@natural parameter} and $A$ the log-partition function.
\end{definition}

\section{Main Theorems and Proofs}

\begin{theorem}[Neyman--Fisher Factorisation \cite{halmos49}]\label{thm:faktorizasyon}\index{factorisation theorem@factorisation theorem}
$T(\X)$ is sufficient for $\theta$ $\iff$ there exist non-negative functions $g$ and $h$ such that for every $\x$ and $\theta$
\[
L_n(\theta; \x) = g(T(\x), \theta)\, h(\x).
\]
\end{theorem}
\begin{proof}[Proof (discrete case)]
($\Rightarrow$) Let $T$ be sufficient and $t=T(\x)$. Then $L(\theta;\x)=\P_\theta(\X=\x)=\P_\theta(\X=\x,T=t)=\P_\theta(T=t)\,\P(\X=\x\mid T=t)$. The second factor is free of $\theta$ by sufficiency; take $g(t,\theta)=\P_\theta(T=t)$ and $h(\x)=\P(\X=\x\mid T=T(\x))$.

($\Leftarrow$) Suppose the factorisation holds. For $\x$ with $T(\x)=t$,
\[
\P_\theta(\X=\x\mid T=t)=\frac{\P_\theta(\X=\x)}{\P_\theta(T=t)}=\frac{g(t,\theta)h(\x)}{\sum_{\mathbf y:T(\mathbf y)=t}g(t,\theta)h(\mathbf y)}=\frac{h(\x)}{\sum_{\mathbf y:T(\mathbf y)=t}h(\mathbf y)},
\]
which is free of $\theta$. (If $T(\x)\neq t$, the conditional probability\index{conditional probability@conditional probability} is zero anyway.) In the continuous case the proof requires measure-theoretic tools (Halmos--Savage).
\end{proof}

\begin{theorem}[Criterion for Minimal Sufficiency \cite{lehmann-scheffe50}]\label{thm:minimal}
If for every pair of sample points $\x,\mathbf y$
\[
\text{the ratio }\frac{L(\theta;\x)}{L(\theta;\mathbf y)}\ \text{is free of }\theta\iff T(\x)=T(\mathbf y),
\]
then $T$ is minimal sufficient\index{minimal sufficient statistic@minimal sufficient statistic}.
\end{theorem}
\begin{proof}[Proof sketch]
\emph{Sufficiency:} For each value $t$ choose a representative $\x_t$ with $T(\x_t)=t$. Since $T(\x)=T(\x_{T(\x)})$, the ratio $h(\x)=L(\theta;\x)/L(\theta;\x_{T(\x)})$ is free of $\theta$; with $g(t,\theta)=L(\theta;\x_t)$ we obtain a factorisation.
\emph{Minimality:} Let $U$ be another sufficient statistic\index{sufficient statistic@sufficient statistic} with $U(\x)=U(\mathbf y)$. By factorisation, $L(\theta;\x)/L(\theta;\mathbf y)=h_U(\x)/h_U(\mathbf y)$ is free of $\theta$; by assumption $T(\x)=T(\mathbf y)$. So $T$ is a function of $U$. (This proof is given for discrete or density models, ignoring ``almost sure'' details; see \cite[Theorem~6.2.13]{casella}, and for the general case \cite[Section~1.6]{lehmann-casella}.)
\end{proof}

\begin{theorem}[Sufficiency in Exponential Families]\label{thm:ustel-yeterli}
If $X_1,\dots,X_n$ are i.i.d.\ from the $k$-parameter exponential family\index{exponential family@exponential family} of Definition~\ref{def:ustel-aile}, then
\[
\mathbf T(\X)=\Big(\sum_{i=1}^n t_1(X_i),\dots,\sum_{i=1}^n t_k(X_i)\Big)
\]
is sufficient. If the set $\{(w_1(\bm\theta),\dots,w_k(\bm\theta)):\bm\theta\in\Theta\}$ contains a non-empty open set in $\R^k$ (``full rank''), $\mathbf T$ is minimal sufficient.
\end{theorem}
\begin{proof}
$L=\prod h(x_i)\cdot c(\bm\theta)^n\exp\{\sum_j w_j(\bm\theta)\sum_i t_j(x_i)\}$; the factorisation can be read off directly. For minimality use Theorem~\ref{thm:minimal}: $L(\x)/L(\mathbf y)=\frac{\prod h(x_i)}{\prod h(y_i)}\exp\{\sum_jw_j(\bm\theta)[T_j(\x)-T_j(\mathbf y)]\}$; under the open-set condition this is free of $\theta$ if and only if $\mathbf T(\x)=\mathbf T(\mathbf y)$.
\end{proof}

\begin{theorem}[Rao--Blackwell Theorem \cite{rao45,blackwell47}]\label{thm:rao-blackwell}\index{rao blackwell theorem@Rao--Blackwell theorem}
Let $W$ be an unbiased estimator\index{unbiased estimator@unbiased estimator} of $g(\theta)$ ($\E_\theta[W] = g(\theta)$, $\Var_\theta W<\infty$) and $T$ a sufficient statistic. Define $W^* = \E[W \mid T]$. Then for every $\theta$:
\begin{enumerate}
  \item $W^*$ is a statistic (i.e.\ it does not depend on $\theta$);
  \item $\E_\theta[W^*] = g(\theta)$;
  \item $\Var_\theta(W^*) \le \Var_\theta(W)$, with equality only if $W$ is already a function of $T$ ($W=W^*$ a.s.).
\end{enumerate}
\end{theorem}
\begin{proof}
(1) Since $T$ is sufficient, the conditional distribution of $\X\mid T$, and hence of $W(\X)\mid T$, does not depend on $\theta$; so $\E[W\mid T]$ is a known function of $T$ alone. \emph{This is where sufficiency is crucial}: conditioning on a non-sufficient $T$ produces an ``estimator'' that depends on $\theta$ and cannot be used.
(2) Tower rule\index{tower property@tower property} (Theorem~\ref{thm:kule}): $\E_\theta[W^*]=\E_\theta[\E(W\mid T)]=\E_\theta W=g(\theta)$.
(3) Law of total variance\index{law of total variance@law of total variance} (Theorem~\ref{thm:toplam-varyans}):
\[
\Var(W) = \Var(\E[W \mid T]) + \E[\Var(W \mid T)] = \Var(W^*) + \E[\Var(W \mid T)] \ge \Var(W^*).
\]
Equality holds when $\E[\Var(W\mid T)]=0$, i.e.\ $W=\E[W\mid T]$ a.s.
\end{proof}

\begin{remark}[Generalisation]
Even without unbiasedness, the conditional Jensen\index{jensen s inequality@Jensen's inequality} inequality gives $R(\theta,W^*)\le R(\theta,W)$ for every convex loss $L(\theta,\cdot)$. So under convex loss it always suffices to restrict attention to estimators that are functions of the sufficient statistic.
\end{remark}

\section{Worked Examples}

\begin{example}[Bernoulli: Sufficiency from the Definition]\label{ex:bernoulli-yeterli}
Let $X_i\iid\text{Bernoulli}(p)$ and $T=\sum X_i\sim\text{Bin}(n,p)$. For every $\x$ with $\sum x_i=t$,
\[
\P_p(\X=\x\mid T=t)=\frac{p^t(1-p)^{n-t}}{\binom nt p^t(1-p)^{n-t}}=\frac{1}{\binom nt},
\]
which is free of $p$: given the number of successes, \emph{which} trials were successes carries no information about $p$.
\end{example}

\begin{example}[Normal: $(\bar X,S^2)$]
For $X_i\iid\Normal(\mu,\sigma^2)$, the identity $\sum(x_i-\mu)^2=\sum(x_i-\bar x)^2+n(\bar x-\mu)^2$ gives
\[
L(\mu,\sigma^2;\x)=(2\pi\sigma^2)^{-n/2}\exp\Big\{-\frac{(n-1)s^2+n(\bar x-\mu)^2}{2\sigma^2}\Big\}.
\]
With $h(\x)=1$, factorisation\index{factorisation theorem@factorisation theorem} shows that $(\bar X,S^2)$ is sufficient; since this is a two-parameter exponential family\index{exponential family@exponential family}, it is minimal sufficient\index{minimal sufficient statistic@minimal sufficient statistic}. If $\sigma^2$ is known, $\bar X$ alone is sufficient.
\end{example}

\begin{example}[$\text{Uniform}(0,\theta)$: the Importance of the Indicator]\label{ex:uniform-yeterli}
\[
L(\theta;\x)=\prod_{i=1}^n\frac1\theta\I\{0<x_i<\theta\}=\theta^{-n}\,\I\{x_{(n)}<\theta\}\cdot\I\{x_{(1)}>0\}.
\]
With $g(t,\theta)=\theta^{-n}\I\{t<\theta\}$ and $h(\x)=\I\{x_{(1)}>0\}$, $T=X_{(n)}$ is sufficient. Forgetting the indicators and writing $L=\theta^{-n}$ is a common mistake that leads to wrong conclusions such as ``nothing is sufficient''.
\end{example}

\begin{example}[Members and Non-Members of the Exponential Family]
\begin{center}
\begin{tabular}{@{}llll@{}}
\toprule
Distribution & Natural parameter $\eta$ & $t(x)$ & Sufficient statistic \\
\midrule
$\text{Bin}(m,p)$ & $\ln\frac{p}{1-p}$ & $x$ & $\sum X_i$ \\
$\text{Poisson}(\lambda)$ & $\ln\lambda$ & $x$ & $\sum X_i$ \\
$\text{Exp}(\lambda)$ & $-\lambda$ & $x$ & $\sum X_i$ \\
$\Normal(\mu,\sigma^2)$ & $(\mu/\sigma^2,\,-1/(2\sigma^2))$ & $(x,x^2)$ & $(\sum X_i,\sum X_i^2)$ \\
$\text{Gamma}(\alpha,\beta)$ & $(\alpha-1,\,-\beta)$ & $(\ln x,x)$ & $(\sum\ln X_i,\sum X_i)$ \\
$\text{Beta}(a,b)$ & $(a-1,\,b-1)$ & $(\ln x,\ln(1-x))$ & $(\sum\ln X_i,\sum\ln(1-X_i))$ \\
\bottomrule
\end{tabular}
\end{center}
Distributions whose support depends on the parameter ($\text{Uniform}(0,\theta)$, the shifted exponential) and location families such as the Cauchy, $t$ and logistic are not exponential families. For example, in the Cauchy location family the minimal sufficient statistic\index{sufficient statistic@sufficient statistic} is the full set of order statistics\index{order statistics@order statistics} $(X_{(1)},\dots,X_{(n)})$: no dimension reduction is possible.
\end{example}

\begin{example}[Rao--Blackwell: Estimating $e^{-\lambda}$ for the Poisson]\label{ex:poisson-rb}
Let $X_i\iid\text{Poisson}(\lambda)$ and $g(\lambda)=\P(X=0)=e^{-\lambda}$ (for example, the probability that no customer arrives in an hour). A crude unbiased estimator\index{unbiased estimator@unbiased estimator} is $W=\I\{X_1=0\}$. $T=\sum X_i$ is sufficient, and $X_1\mid T=t\sim\text{Bin}(t,1/n)$ (the conditional distribution of independent Poissons). Hence
\[
W^*=\E[W\mid T]=\P(X_1=0\mid T)=\Big(1-\frac1n\Big)^{T}.
\]
Using the probability generating function of $T\sim\text{Poisson}(n\lambda)$, $\E[(1-1/n)^T]=e^{n\lambda(-1/n)}=e^{-\lambda}$ (unbiased), and
\[
\Var(W^*)=e^{-2\lambda}\big(e^{\lambda/n}-1\big)\quad\ll\quad\Var(W)=e^{-\lambda}(1-e^{-\lambda}).
\]
The intuitive ``plug-in\index{plug in principle@plug-in principle}'' estimator $e^{-\bar X}$ is biased, whereas $W^*$ is unbiased and will be shown to be the UMVUE\index{umvue@UMVUE} in Chapter~\ref{ch:tamlik}.
\end{example}

\begin{example}[Rao--Blackwell: $\text{Uniform}(0,\theta)$]
$W=2X_1$ is unbiased. Given $X_{(n)}=m$, $X_1$ equals $m$ with probability $1/n$ and is otherwise $\text{Uniform}(0,m)$. Hence
\[
\E[2X_1\mid X_{(n)}=m]=2\Big[\frac mn+\frac{n-1}{n}\cdot\frac m2\Big]=\frac{n+1}{n}m .
\]
The Rao--Blackwell\index{rao blackwell theorem@Rao--Blackwell theorem} procedure mechanically produces the ``intuitive'' estimator $\frac{n+1}{n}X_{(n)}$ of Chapter~\ref{ch:nokta-tahmin}.
\end{example}

\section{Computational Applications}

\subsection{Bernoulli: $X_1\to\bar X$}

\pycode{code/ch09/c09_1_bernoulli.py}
\begin{lstlisting}[style=output]
Var(W=X_1): 0.2400 | Var(W*=X_bar): 0.0240
Theory    : 0.2400                | 0.0240
\end{lstlisting}

$\E[X_1\mid\sum X_i=t]=t/n$ (by symmetry), i.e.\ $W^*=\bar X$ and the variance falls by a factor of $n$.

\subsection{Poisson: Three Estimators of $\P(X=0)$}

\pycode{code/ch09/c09_2_poisson_three_estimators.py}
\begin{lstlisting}[style=output]
Target e^-lambda = 0.13534 | CRLB = 3.663e-03
W  = I(X1=0)         | mean=0.13553 | bias=+0.00020 | Var=1.172e-01 | MSE=1.172e-01
W* = (1-1/n)^T (RB)  | mean=0.13539 | bias=+0.00005 | Var=4.042e-03 | MSE=4.042e-03
exp(-Xbar) (plug-in) | mean=0.14914 | bias=+0.01380 | Var=4.399e-03 | MSE=4.590e-03
Theory Var(W*) = 4.055e-03
\end{lstlisting}

The Rao--Blackwell\index{rao blackwell theorem@Rao--Blackwell theorem} estimator reduces the variance by a factor of about $29$. Its variance is slightly above the CRLB\index{cramer rao lower bound crlb@Cramér--Rao lower bound (CRLB)}: since $e^{-\lambda}$ is not a linear function of $\E T$, no efficient estimator\index{efficient estimator@efficient estimator} exists (Corollary~\ref{cor:crlb-esitlik}). The intuitive plug-in\index{plug in principle@plug-in principle} estimator $e^{-\bar X}$ is biased upwards by Jensen\index{jensen s inequality@Jensen's inequality}'s inequality ($e^{-x}$ is convex), and in this example both its variance and its MSE exceed those of $W^*$: here the Rao--Blackwell procedure beats intuition not only in unbiasedness but also in MSE. (This need not always be so; in Example~\ref{ex:en-iyi-bolen} the biased estimator was better.)

\subsection{$\text{Uniform}(0,\theta)$: Rao--Blackwell by Simulation}

\pycode{code/ch09/c09_3_rao_blackwell_simulation.py}
\begin{lstlisting}[style=output]
M slice 0: mean(W | M) = 4.078  vs  (n+1)/n * mean(M) = 4.089
M slice 1: mean(W | M) = 4.822  vs  (n+1)/n * mean(M) = 4.827
M slice 2: mean(W | M) = 5.143  vs  (n+1)/n * mean(M) = 5.153
M slice 3: mean(W | M) = 5.381  vs  (n+1)/n * mean(M) = 5.377
M slice 4: mean(W | M) = 5.553  vs  (n+1)/n * mean(M) = 5.550
Var(W) = 8.3277 (theory 8.3333) | Var(W*) = 0.3123 (theory 0.3125)
\end{lstlisting}

\section{Comments and Notes}

\begin{remark}[Sufficiency Depends on the Model]
$\bar X$ is sufficient for $\mu$ in the normal model ($\sigma^2$ known), but not in the Laplace or Cauchy location models. If the model is wrong, reducing to the ``sufficient'' statistic can lose information; that is why the raw data are always needed for model checking (residual analysis).
\end{remark}

\begin{remark}[The Pitman--Koopman--Darmois Theorem]
Among regular families whose support does not depend on the parameter, only exponential families\index{exponential family@exponential family} admit a sufficient statistic\index{sufficient statistic@sufficient statistic} whose dimension does not grow with $n$. That is, ``lossless compression of the data into a fixed-dimensional summary'' is essentially a property of exponential families. This explains their central place in statistics (GLMs\index{generalised linear model glm@generalised linear model (GLM)}, conjugate\index{conjugate prior@conjugate prior} priors).
\end{remark}

\begin{remark}[Rao--Blackwell in Practice]
The Rao--Blackwell\index{rao blackwell theorem@Rao--Blackwell theorem} principle is also used in Monte Carlo: when estimating an expectation by simulation, integrating out analytically the part whose conditional expectation\index{conditional expectation@conditional expectation} can be computed reduces the variance (``Rao--Blackwellised'' MCMC\index{mcmc@MCMC} estimators, conditional Monte Carlo).
\end{remark}

% --- CHAPTER SUMMARY ---
\section{Chapter Summary}
\begin{summary}
\begin{itemize}
  \item If $T$ is sufficient, the conditional distribution of the data given $T$ is free of $\theta$ (Definition~\ref{def:yeterli}). $T$ carries all the information in the data about $\theta$.
  \item The factorisation theorem is the practical criterion for sufficiency: $f(\x;\theta)=g(T(\x);\theta)h(\x)$ (Theorem~\ref{thm:faktorizasyon}).
  \item A minimal sufficient statistic is the coarsest sufficient reduction. It is characterised by likelihood ratios being free of $\theta$ (Theorem~\ref{thm:minimal}).
  \item In an exponential family, $\big(\sum_it_1(X_i),\dots,\sum_it_k(X_i)\big)$ is sufficient (Theorem~\ref{thm:ustel-yeterli}); its dimension does not depend on $n$.
  \item Rao--Blackwell theorem (Theorem~\ref{thm:rao-blackwell}): conditioning any estimator on a sufficient statistic preserves its bias and does not increase its MSE. Good estimators should be sought among functions of the sufficient statistic.
\end{itemize}
\end{summary}
% --- END OF CHAPTER SUMMARY ---

\paragraph{Further reading.} \cite[Chapter~6]{casella}, \cite[Section~1.6]{lehmann-casella}, \cite[Chapter~2]{shao}.

\section{Exercises}

\begin{exercise}\label{exr:09-1}
Let $X_i\iid\text{Gamma}(\alpha,\beta)$ with both parameters unknown. Find a minimal sufficient statistic. What changes if $\alpha$ is known?
\hfill\hyperref[coz:09-1]{\footnotesize$\rightarrow$\,Solution}
\end{exercise}

\begin{exercise}\label{exr:09-2}
For $X_i\iid\text{Uniform}(\theta,\theta+1)$, show using Theorem~\ref{thm:minimal} that $(X_{(1)},X_{(n)})$ is minimal sufficient. What does the existence of a two-dimensional minimal sufficient statistic in a one-parameter model indicate? (See Chapter~\ref{ch:tamlik}.)
\hfill\hyperref[coz:09-2]{\footnotesize$\rightarrow$\,Solution}
\end{exercise}

\begin{exercise}\label{exr:09-3}
Let $X_i\iid\text{Bernoulli}(p)$ and $g(p)=p^2$. $W=X_1X_2$ is unbiased. Compute $W^*=\E[W\mid T]$ ($T=\sum X_i$) and show that $W^*=\frac{T(T-1)}{n(n-1)}$.
\hfill\hyperref[coz:09-3]{\footnotesize$\rightarrow$\,Solution}
\end{exercise}

\begin{exercise}\label{exr:09-4}
Let $X_i\iid\Normal(\mu,1)$ and $g(\mu)=\P(X\le c)=\Phi(c-\mu)$. Rao--Blackwellise $W=\I\{X_1\le c\}$ with respect to $\bar X$. (Hint: $X_1\mid\bar X\sim\Normal(\bar X,1-1/n)$.) Compare the result with the plug-in estimator $\Phi(c-\bar X)$ by simulation.
\hfill\hyperref[coz:09-4]{\footnotesize$\rightarrow$\,Solution}
\end{exercise}

\end{document}
